% Text of the short version (paper/short.tex). Every statement is one of the paper
% (sections/01 to 09) and every number of a run is the same \result{} key; the predictions and
% the finite-element numbers are typed as the paper types them (sections/05-results.tex,
% sections/fem-reference.tex). It cites the paper and the supplement by section number
% (Sections 5.4 and 7 of the paper, S1 to S8): change them if that numbering changes. If the
% paper's findings change, change this file with them.

\section{Introduction}\label{sh:introduction}

Contact between the bodies of a multibody system is most often modelled at a point: a penalty
force that grows with the penetration depth \cite{HuntCrossley1975}, or a unilateral
constraint with Coulomb friction. Both assume that the contact patch is small against the bodies.
Many contacts of practical interest are \emph{conforming} instead: a flat part on a table, a
pin in a hole, a journal in a bearing. There the friction torque of a spinning disk or the
tilt of an eccentric load depend on how the pressure is distributed over the patch.

Several families of models address the patch. Elastic-foundation models cover a body with a
bed of independent springs and sum their forces over a discretized surface,
as the polygonal contact model does \cite{Hippmann2004}. Volumetric models
\cite{Gonthier2005,Gonthier2007thesis} integrate the same foundation in closed form for a flat
patch and express the wrench through the volume, the centroid and the inertia tensor of the
overlap of the undeformed bodies. Pressure-field
(hydroelastic) contact \cite{Elandt2019}, the patch model of Drake \cite{Drake},
defines a pressure field inside each body and a contact surface where the two fields are equal.
These models have been validated on a few cases each; the comparisons of
simulators from robotics test point contacts \cite{Erez2015}, and the benchmarks of the
multibody community are built on mechanisms with ideal joints \cite{Gonzalez2006}. We did
not find a set of conforming-contact problems with analytic references that has been run
across model families and across engines.

This paper offers such a set and uses it. Its origin is the second author's doctoral thesis
\cite{Vilela2018}, which evaluated Gonthier's volumetric model on the exact intersection of two
triangle meshes inside one multibody code. The work reported here re-implements that pipeline
as an engine-independent library, adds a second law that integrates the same foundation over
the contact surface instead of lumping it, and runs both, next to baselines from the
literature and the native contact models of Drake, Project Chrono \cite{Tasora2016} and MuJoCo
\cite{Todorov2012}, on a battery of problems with known answers. The paper is a benchmark and
a comparison; it does not propose a new contact model. This short version gives the laws, the
benchmark and the main findings; the full paper and its supplementary material, cited as
Sections~S1 to~S8, give the rest (Section~\ref{sh:conclusions}).

\section{The two laws on exact mesh intersections}\label{sh:laws}

Both laws are the Winkler foundation applied to the overlap of two undeformed bodies $A$ and
$B$ given by closed triangle meshes. The overlap is read as a bundle of columns along a normal
$\n$; a column of length $\ell$ stands for a spring compressed by $\ell$ and presses with
$K\ell$ per unit of area, with the stiffness $K$ in N/m$^3$. A rate term of Hunt--Crossley form
multiplies the pressure by $1 - h\,v_n$, with $v_n = \n\cdot\vv$ the rate at which the bodies
separate and $h$ [s/m] the damping factor, and friction is Coulomb's, with coefficient $\mu$;
$\vv$ and $\w$ are the velocity field and the angular velocity of $A$ relative to $B$.

The boundary of the overlap is obtained by clipping the triangles of each mesh against the
other, and the divergence theorem gives its volume $V$, its centroid $\vect{c}$ and its inertia
tensor $\tens{J}$ (unit density, about $\vect{c}$), exact for the polyhedra given up to
round-off. The normal $\n$ is the direction
of the vector area of the part of the boundary that belongs to $B$. Both laws use one normal
per contact region, and the distributed law requires convex shapes, so a non-convex body is
given as convex pieces that fill it exactly.

\paragraph{Gonthier's volumetric law.} The law \law{gonthier} implements the model of
Gonthier's thesis \cite{Gonthier2007thesis} on the exact $V$, $\vect{c}$ and $\tens{J}$. The
normal force acts along $\n$ at the centroid, with a rolling-resistance torque, and a Coulomb
force and a Coulomb spinning torque share one friction budget:
\begin{equation}
  \Fn = K V \left(1 - h\,v_n\right), \;\;
  \vect{M}_r = -K h\, \tens{J}\, \w_t , \;\;
  \vect{f}_t = -\mu\,\Fn\,\frac{\vv_t}{v_\mathrm{avg}}, \;\;
  \vect{M}_s = -\mu\,\Fn\, r_g^2\,\frac{\omega_n}{v_\mathrm{avg}}\,\n ,
  \label{sh:gonthier}
\end{equation}
with $v_n$ and the tangential velocity $\vv_t$ taken at the centroid, $\w_t$ and
$\omega_n$ the tangential and normal parts of $\w$, $r_g^2 = \n\cdot\tens{J}\,\n/V$ the squared
radius of gyration of the overlap about the normal and
$v_\mathrm{avg}^2 = |\vv_t|^2 + r_g^2\,\omega_n^2$. The lumping has a known price: for a flat
circular patch of radius $R$, $r_g = R/\sqrt2$, whereas the exact Coulomb integral of a uniform
pressure has the arm $2R/3$, so the spinning torque of the published model is 1.06 times the
exact one.

\paragraph{The distributed law.} The law \law{distributed} evaluates the same foundation point
by point on the contact surface. Through a point $\x$ of a face of either body a column runs
along $\pm\n$ to the first face of the surface of either body, and the pressure at $\x$, along
the normal of its face, is
\begin{equation}
  p(\x) = \tfrac12\,K\,\ell(\x)\left(1 - h\,\n\cdot\vv(\x)\right),
  \label{sh:pressure}
\end{equation}
each end of a column carrying half, so that the law is symmetric in $A$ and $B$; where the
damping factor is negative there is no pressure, since there is no adhesion. The column length
is piecewise affine, so the normal wrench is integrated exactly for the polyhedra, with no
quadrature error and no dependence on the triangulation: one triangle per flat face is enough
(Section~S1). For a flat contact the law reduces to the normal force of
Eq.~\eqref{sh:gonthier}, and the first moment of the damped pressure is, to first order in the
depth, the rolling resistance $\vect{M}_r$: under a rolling body the leading half of the patch
is loaded more than the trailing half.

\paragraph{Coulomb friction in closed form.} Every triangle of the contact surface carries a
linear pressure $p$ and slips with the relative velocity field of the two bodies. The friction
traction $-\mu\,p\,\vv_t/|\vv_t|$ changes direction over the triangle, so its integral is not a
polynomial; it can nevertheless be done in closed form, because the slip field of a rigid
motion in a plane is a rotation about a point $\q$, the instantaneous centre of the slip. With
$\m$ the unit normal of the triangle $T$, $\rr = \x - \q$, $\sigma$ the sign of $\w\cdot\m$ and
the pressure written as $p = p_q + \nabla p\cdot\rr$, the friction force and its moment about
$\q$ are
\begin{equation}
  \vect{F}_f = -\mu\,\sigma\,\m\times\left(p_q\,\vect{U} + \tens{W}\,\nabla p\right), \qquad
  \vect{M}_f = -\mu\,\sigma\left(p_q\,D + \nabla p\cdot\vect{R}\right)\m ,
  \label{sh:friction}
\end{equation}
with four polar moments of the triangle,
$D = \int_T |\rr|\,\dd A$, $\vect{R} = \int_T \rr\,|\rr|\,\dd A$,
$\vect{U} = \int_T \rr/|\rr|\,\dd A$ and $\tens{W} = \int_T \rr\,\rr^{\mathsf T}/|\rr|\,\dd A$.
$D$ and $\vect{U}$ are the moments that Jia \cite{Jia1997} gives for a polygon under uniform
pressure; $\vect{R}$ and $\tens{W}$ are what the linear term adds. All four are evaluated edge
by edge with an inverse hyperbolic sine and powers of a square root, with no case distinction
for $\q$ inside, outside or on the boundary of the triangle (Section~S6 gives the formulas and
their derivation). Earlier work on distributed contact interpolates the friction from the
vertices or samples it at one point per face \cite{Hasegawa2004,Hippmann2004}. This closed
form is the only element of \law{distributed} that we claim as new; the rest of the law
belongs to the line that goes from Hasegawa and Sato \cite{Hasegawa2004} to hydroelastic
contact \cite{Elandt2019}.

\paragraph{Stiction and coupling.} Both laws hold a contact at rest with a bristle, a
spring--damper on the tangential deflection whose stiffness is proportional to the load. The backends for Drake, Project Chrono
and MuJoCo switch the engine's own contact off and apply the wrenches of the library as
external forces, computed at the start of the step and held during it. This explicit coupling
has a stability limit: for the normal stiffness of a flat contact of area $A_c$, with $w$ the
inverse masses that the contacts of a body see and $\dt$ the step,
\begin{equation}
  K A_c\,\dt^2 \sum w < 4 ,
  \label{sh:stability}
\end{equation}
and $<2$ with the dampers at their bounds, so a body of mass $m$ held between two flat contacts
needs $\dt < \sqrt{m/(K A_c)}$. An implicit coupling needs the derivatives of the wrench, which
the laws do not provide yet.

\paragraph{Baselines.} Two variants of \law{distributed} replace only the closed form by the
quadratures of the literature: the vertex rule of Hasegawa and Sato \cite{Hasegawa2004} and a
one-point rule at the pressure centroid of each triangle. The law \law{pcm} is a per-face
elastic foundation after the polygonal contact model \cite{Hippmann2004}: one spring per
triangle, with a depth measured along local normals and a friction without a stick state.

\section{The benchmark}\label{sh:benchmark}

The battery has \result{battery/scenarios} scenarios, each with a set-up, a reference solution
and an error measure, written once against an engine-independent interface
(Table~\ref{sh:tab:scenarios}). Every scenario is classified by what its reference rests on.
\emph{Class A}: a rigid-body result that does not depend on the contact model, such as the
acceleration of a block on an incline. Only these
scenarios compare a compliant law with a contact model of another kind on equal terms.
\emph{Class B}: a reference that assumes a pressure distribution, such as the stop time of a
spinning disk. The nominal reference uses the distribution that a rigid body produces on a
Winkler foundation, so agreement with it is in part agreement with the model under test; other
distributions are defensible, and the part of the reference that does not depend on the model is the \emph{band}
that they span. For the disk, the friction torque over $\mu N R$, with $N$ the load, is $2/3$
for uniform pressure, $1/2$ for uniform wear and $\pi/4$ for the rigid punch. A method inside
the band is not wrong. \emph{Class C}: consistency with the Winkler foundation; the reference
contains $K$ or integrates the law's own force. These scenarios say whether a method is a
Winkler foundation of stiffness $K$, not whether it is accurate.

\begin{table}[t]
\caption{The scenarios, by class (Section~S7 gives the parameters and the reference of each)}
\label{sh:tab:scenarios}
\footnotesize
\begin{tabular}{@{}p{0.05\linewidth}p{0.40\linewidth}p{0.48\linewidth}@{}}
\toprule
Class & Reference & Scenarios \\
\midrule
A & Rigid-body results: sliding and rolling on planes, in a bowl and on a turntable, impacts,
    bodies that must be held, a ground cut into pieces, a shaft in a hole with clearance &
  \scn{block}, \scn{cylinder}, \scn{sphere}, \scn{push}, \scn{billiard}, \scn{vgroove},
  \scn{twoblocks}, \scn{bowl}, \scn{trough}, \scn{journal}, \scn{pieces1}, \scn{2}, \scn{8},
  \scn{64}, \scn{wedge}, \scn{oblique}, \scn{edge}, \scn{turntable}, \scn{energy}, \scn{tower},
  \scn{massratio} \\
B & Stop time of a spinning body under an assumed pressure, with its band &
  \scn{disk}, \scn{contensou}, \scn{ring}, \scn{bearing} (a shaft in a hole of its own radius),
  \scn{cone} (a frustum in a conical seat) \\
C & The foundation's own predictions: tipping angle, rebound, depth of a press fit, tilt &
  \scn{tip}, \scn{drop}, \scn{hole}, \scn{corner}, \scn{offset}, \scn{stack},
  \scn{restitution} \\
\botrule
\end{tabular}
\end{table}

The error of a run is the worst normalized deviation of every quantity that the reference
fixes. Errors are not commensurable across scenarios, so methods are
compared within a scenario. Parameters are fixed before any run; the
foundations have $K$ between $10^6$ and $10^8$\,N/m$^3$ and $h = 5$\,s/m. Every scenario is run
in a base variant (curved surfaces with 64 segments per turn, $\dt = 1$\,ms) and in series of
the step (2 to 0.1\,ms), the mesh (32 and 128 segments), the stiffness ($K\times10$) and
perturbed initial velocities, with the methods of the library on each of the three engines and
with six native models, in \result{battery/runs} runs.

The native models stand for three families. \emph{Point contact}: Drake's compliant point
contact, Chrono's penalty model (SMC) and MuJoCo's soft constraints with two sets of
parameters; a contact between two faces is a few points with a stiffness in N/m.
\emph{Complementarity contact}: Chrono NSC, rigid contact with Coulomb friction.
\emph{Pressure-field contact}: Drake hydroelastic, configured so that its pressure is $K$
times the depth; it is then the same pressure model as the two laws near the surface, in
another discretization. The configurations are ours and were not reviewed by the developers of
the engines, so we draw from the native runs only what follows from a property of the model
family (Section~S3 gives the configurations and every other result).

\section{Findings}\label{sh:findings}

Each law gives the same errors on the three engines, within \result{independence/worst}
wherever a run holds (a run holds when it finishes with an error below 100\,\%, and is lost
when the error is larger); the three backends share no code, so this is a check of the
implementation. Their step series are mostly flat (\result{convn/dt/distributed/flat} of
\result{convn/dt/distributed/n} for \law{distributed}), so the errors at the base step are
mostly those of the model and the mesh, not of the time integration.

\subsection{Class A: where the laws fail, and why}\label{sh:classA}

Where a body slides or rolls on one plane both laws are within 1\,\%, except for the time at
which rolling starts in \scn{billiard} (\result{err/billiard/gonthier} and
\result{err/billiard/distributed}), and they hold the bodies that must be held
(\scn{push}, \scn{wedge}, \scn{tower}). The methods without a stick state creep long before the friction limit (\law{pcm}:
\result{err/push/pcm} in \scn{push}). The larger errors of the two laws have causes that can
be predicted.

\emph{Rolling resistance of the damped foundation.} In \scn{turntable}, a ball rolling on a
free rotating table whose rigid orbit is a closed circle, the orbit rate is right and the
radius grows by \result{err/turntable/distributed}. The cause is the torque $\vect{M}_r$ of
Eq.~\eqref{sh:gonthier}: the no-slip equations with it predict an orbit that grows by 77\,\%
over the run.
% 77 %: a prediction, typed as in sections/05-results.tex (docs/history.md).
This is the model's answer at the $h$ and $K$ of the scenario, not a numerical error: any
damped Winkler foundation has this torque, and with $K\times10$ the error falls to \result{err/turntable/distributed/stiff}.

\emph{The stability bound of the explicit coupling.} In \scn{massratio} a 1\,kg block lies
between a 1000\,kg block and the ground, and Eq.~\eqref{sh:stability} requires a step below
0.25\,ms. The two laws lose the block at 2, 1 and 0.5\,ms and hold it from 0.2\,ms, with
\result{err/massratio/distributed/dt0.2} of error there. Complementarity contact, implicit in
the contact force, holds the scenario at the base step.

\emph{One normal per pair of convex pieces.} In \scn{journal}, a shaft spinning in a hole with
clearance, the error of \law{distributed} does not fall with the mesh:
\result{err/journal/distributed/mesh32}, \result{err/journal/distributed} and
\result{err/journal/distributed/mesh128} at 32, 64 and 128 segments per turn, while \law{pcm},
whose depth follows the local normal, reaches \result{err/journal/pcm/mesh128} at 128. The
geometry of \law{distributed} is exact for the polyhedra, but its answer for a curved
conforming contact does not converge to that of the curved bodies.

\subsection{Class B: the band, and a finite-element reference}\label{sh:classB}

Table~\ref{sh:tab:errors-B} gives the errors against the nominal reference and marks the
position of each run against the band. No elastic-foundation method is outside the band in any
scenario: \law{distributed}, \law{gonthier}, \law{pcm} and Drake hydroelastic have no run
marked as outside (the variant with the friction of the 2018 thesis is the exception among the
methods of the library). \law{distributed} and Drake hydroelastic, the two methods that integrate the
foundation over the contact surface, also meet the nominal value, within
\result{err/bearing/distributed} and \result{err/cone/hydro}. That agreement is a consistency
result: the nominal value is the model's own distribution, so two independent implementations
of the same pressure model integrate it correctly, and nothing follows about whether that
pressure is the true one. \law{gonthier} misses the nominal value by what its lumping
predicts, \result{err/disk/gonthier} in \scn{disk}. A point-contact model does not represent
the distribution: a face held by points of its rim has the friction torque of a ring,
$\mu N R$ against $\tfrac23\mu N R$ for the disk, unless the user supplies a torsional
friction coefficient for each geometry; the same holds for complementarity contact on the same
contact points. Each of these models is outside the band in
\result{band/smc/outside} to \result{band/dpoint/outside} of the five scenarios.

\begin{table}[t]
\caption{Class B: error [\%] against the nominal reference, base variant. Gon: \law{gonthier};
G18: its variant with the friction of the 2018 thesis; Dis: \law{distributed}; Dr-h, Dr-p: Drake hydroelastic and point; Ch-S, Ch-N:
Chrono SMC and NSC; Mj-s, Mj-d: MuJoCo soft and default. Band: in percent of the nominal stop
time. $^\dagger$: outside the band; $^{*}$: undetermined (the sign of the error was not
recorded); unmarked: inside; \textit{nan}: no error value. Native models: results of our
configurations}
\label{sh:tab:errors-B}
\generatedtable{errors_B}
\end{table}

To place a continuum in the band we solved by finite elements the static frictionless contact
of a linear-elastic support with a rigid body, for \scn{disk} and \scn{bearing}. A thin layer
is a Winkler foundation, with a nearly uniform pressure under the disk. The slab of the
battery is thick, and there the shear between neighbouring fibres, which a Winkler foundation
omits, moves the load to the rim: a continuum support puts \scn{disk} at the rigid-punch end
of its band, with a stop time 15\,\% shorter than the nominal one. Against that reference
\law{distributed} and Drake hydroelastic, which reproduce the uniform pressure of their own
model, stop 18\,\% late. This is a limit of the Winkler foundation itself on a thick support,
shared by its two independent integrations; a finer mesh or a smaller step would not reduce
it. For \scn{bearing} the continuum sleeve gives, for $\nu=0.3$, the cosine pressure of the foundation, the
four foundation methods are within 5\,\% of the finite-element stop time and the reference
cannot rank them (Section~5.4 of the full paper).
% 15 %, 18 %, 5 %: typed from sections/fem-reference.tex (tools/fem, results/fem).

On class C, a consistency check, \law{distributed} is within 5\,\% in
\result{pass/C/distributed/5} scenarios and \law{gonthier} in \result{pass/C/gonthier/5}; on a
ground cut into 64 pieces \law{distributed} gives the same error
(\result{err/pieces64/distributed}).

\subsection{The friction quadrature and the triangulation}\label{sh:quadrature}

On the meshes of the battery the closed-form friction and the two quadratures of the
literature differ widely on the flat contacts that spin: in \scn{disk},
\result{err/disk/distributed} with the closed form against \result{err/disk/vertex} with the
vertex rule and \result{err/disk/centroid} with one point per triangle. The size of those differences is a property of the battery's
meshes, in which the flat cap of a disk is a fan of long triangles from a vertex of its rim.
To separate the rule from the mesh, the three rules were evaluated apart from the library on a
disk spinning under uniform pressure, for several triangulations of the same face
(Fig.~\ref{sh:fig:triangulation}). The closed form gives the same wrench on every
triangulation, to round-off. On the fan from the rim the error of the spin torque is
\result{tri/disk/spin/rim-fan/64/centroid/moment} with the one-point rule and
\result{tri/disk/spin/rim-fan/64/vertex/moment} with the vertex rule, at any number of segments; on well-shaped
triangles both converge at second order. So the closed form is not more accurate than a
well-used quadrature. What it gives is independence: the
friction wrench of a flat face does not depend on how the face is triangulated or on how the
bodies are cut into pieces, so a flat face needs no mesh of its own.

\begin{figure}[t]
\centering
\generatedfigure[width=0.85\linewidth]{ablation_triangulation}
\caption{Error of the spin torque of a disk under uniform pressure with the one-point rule and
with the vertex rule, against the number of triangles of the face, for five families of
triangulations. The closed form has no error on any of them}
\label{sh:fig:triangulation}
\end{figure}

\subsection{A stiffer foundation: friction-induced flutter}\label{sh:flutter}

A foundation ten times stiffer reduces the class A errors that come from the compliance and
the damping, and makes the class B errors of the two laws larger in \result{stiffn/distributed/worse/B} of
the five scenarios: \scn{disk} goes from
\result{err/disk/distributed/dt0.2} to \result{err/disk/distributed/stiff} for
\law{distributed}, at the same step. The cause is not the explicit coupling, since in
diagnostic runs the error was the same at steps from 1 to 0.05\,ms; it is a self-excited
wobble of the model. A tilt of the spinning disk shifts the pressure, so the friction of the
heavier side no longer cancels that of the lighter one, and the net friction force acts at the
contact, half the thickness $H$ of the disk below the centre of mass. Its moment turns the
tilt instead of restoring it, and the whirl grows when
\begin{equation}
  \mu K\,\frac{\pi R^3}{3}\,\frac{H}{2} \;>\; c_\theta\,\omega_\theta ,
  \qquad c_\theta = \frac{h\,m g R^2}{4}, \qquad \omega_\theta = \sqrt{\frac{K I_A}{I_t}} ,
  \label{sh:eq:flutter}
\end{equation}
with $c_\theta$ the rocking damping of the foundation under the weight $m g$,
$I_A = \pi R^4/4$ and $I_t$ the transverse moment of inertia of the disk. For the disk of the scenario the threshold is about twice the scenario's stiffness:
at $K$ the wobble decays, at $3K$ the condition predicts a growth of 6\,s$^{-1}$ and the
diagnostic run showed 6.6, and at $10K$ it predicts 31\,s$^{-1}$ against 32 measured. The disk
then rides on its rim and stops early.
% Typed as in sections/05-results.tex (docs/research/anomalies.md §1).
The left side grows with $K$ and the right side with $\sqrt K$: a foundation made stiffer at
the same $h$ is less damped in rocking. Eq.~\eqref{sh:eq:flutter} is a limit of validity of
the damped Winkler law with sliding friction under a body whose centre of mass is above the
contact, and not a defect of an implementation.

\subsection{Cost and paired comparisons}\label{sh:cost}

In wall-clock time over simulated time, on one processor and at the base step, \law{gonthier}
has a median of \result{cost/gonthier/median}, \law{distributed} of
\result{cost/distributed/median} and Drake hydroelastic, as configured here, of
\result{cost/hydro/median}. Across scenarios, two methods are compared by the geometric mean
of the ratio of their errors on the same scenarios, with a bootstrap interval and Holm's
correction over \result{primary/n} primary comparisons (Section~S5). Only
\result{primary/significant} survive the correction, both on classes A and B together:
\law{distributed} has about half the error of \law{gonthier}
(\result{robust/distributed-vs-gonthier/A+B/full}) and about a quarter of that of Chrono SMC
(\result{robust/distributed-vs-smc/A+B/full}), the latter narrowly. On class A alone nothing
separates the laws from the native models.

\section{Which law to use when}\label{sh:which}

Table~\ref{sh:tab:decision} sets the situations of the battery
against the property of a contact model that each one needs. Hydroelastic contact, the same
pressure model as \law{distributed} near the surface with a curved contact surface and an
implicit solver, serves better a user who needs non-convex geometry, stiff contact or large
steps, and needs a pressure field in every compliant body.

\begin{table}[t]
\caption{Which property of a contact model each situation needs, as far as classes A and B of
the battery show}
\label{sh:tab:decision}
\footnotesize
\begin{tabular}{@{}p{0.36\linewidth}p{0.58\linewidth}@{}}
\toprule
Situation & What it needs \\
\midrule
A rigid body sliding or rolling on one plane & Nothing beyond point or complementarity
  contact; a patch model adds nothing \\
A flat face that spins, or slides and spins & A pressure distributed over the face:
  \law{distributed} or a pressure field; \law{gonthier} where its lumping error is acceptable \\
Curved conforming contact (a journal, a conical seat) & Local normals: a per-face foundation
  or a pressure field. A finer mesh is not a remedy for \law{distributed} \\
Bodies or grounds cut into convex pieces & A wrench independent of the cut: \law{distributed} \\
A body held by friction near its limit & A stick state, or a complementarity formulation \\
Stiff contacts on light bodies, large mass ratios & A formulation implicit in the contact
  force; not the explicit laws beyond Eq.~\eqref{sh:stability} \\
Long rolling & No damping torque in the patch: rigid contact, or the laws with little damping
  or a stiff foundation \\
Only surface meshes are available & \law{gonthier} (any closed mesh) or \law{distributed}
  (convex pieces) \\
\botrule
\end{tabular}
\end{table}

\section{Limitations, conclusions and where to find the rest}\label{sh:conclusions}

Every reference of the battery is analytic, and the finite-element one is static and covers
two scenarios; there is no experimental reference, and nothing here shows that a Winkler
foundation of stiffness $K$ describes a physical contact beyond a thin layer. The foundations are soft, chosen so that the explicit coupling is stable at the base
step, and the damping factor has no physical justification. The scenarios
are few and hand-picked, and the laws were changed after seeing every one of them (the last
\result{battery/heldout} scenarios, written with their references before any law was run on
them, still saw three general changes after their first run), so no part of the
battery is an out-of-sample test of the two laws and the statistics are descriptive. The
results of the native models are results of our configurations and should not be read as the
performance of those engines in expert hands.
Section~7 of the full paper states these limits in full.

The battery, with its classes and its band of admissible pressure distributions, is the
reusable product of the work; the laws are an open implementation of known models. The findings that we expect to survive another battery are those tied to a
mechanism: the rim torque of a point contact, the 6\,\% of Gonthier's lumping, the distance
between the foundation and a continuum on a thick support, the rolling resistance of the
damped foundation, the flutter of Eq.~\eqref{sh:eq:flutter}, the stability bound of
Eq.~\eqref{sh:stability}, and one normal per pair of convex pieces on curved conforming
contact. An implicit coupling, with the Jacobians of the contact
wrench, would remove the step limit, and an experimental reference would turn consistency with
an assumed pressure into validation.

Everything else is at \code{https://dvilelaf.github.io/volumetric-contact/}: the full paper
(\code{paper.pdf}: the related work, the laws in full, the tables of every class, the
discussion), its supplementary material (\code{supplement.pdf}, Sections~S1 to~S8: the
constructions, the step bounds, the native configurations, the complete tables, the statistics,
the derivation of the polar moments, the parameters of every scenario, the wider related work) and a viewer that replays every run. The library, the scenarios, the results and the sources of the papers, with the
scripts that generate every number in them, are open source (BSD 3-Clause) at
\code{https://github.com/dvilelaf/volumetric-contact}.

\backmatter

\section*{Declarations}

\bmhead{Use of generative AI}
This work was carried out by a generative AI system, listed as first author, under human
direction: the software, the design of the benchmark, the analysis of its results and this
manuscript were produced by the model Claude Fable~5.2 (Anthropic), which also directed
further instances of Claude models. The
second author set the questions, decided the scope, and approved the changes to the contact
laws and the rules of the benchmark; the full paper describes the division
of work in full. The first author is an AI system and cannot be held accountable; the second
author directed the work and is the person to whom questions about it should be addressed.

\bmhead{Funding, competing interests, data and code}
No funding was received. The second author has no competing interests to declare; the first
author is a commercial AI system developed by Anthropic. The code and the results are at the address given above.
